% --- Archon physics formalization source begin ---
% archon:physics
% archon:covers PhyXMiniProblems/problem_phyx_mini_0511.lean
% archon:source-report reports/phyx_mini/problem_phyx_mini_0511.source.json

\chapter{Physics problem phyx\_mini\_0511}
\label{ch:PhyXMiniProblems_problem_phyx_mini_0511}

\paragraph{Problem source.}
\#\# Physical scenario

Two protons (each with mass \textbackslash{}( M\_\{\textbackslash{}text\{p\}\} = 1.67 \textbackslash{}times 10\textasciicircum{}\{-27\} \textbackslash{}text\{ kg\} \textbackslash{})) are initially moving with equal speeds in opposite directions. They continue to exist after a head-on collision that also produces a neutral pion of mass \textbackslash{}( M\_\{\textbackslash{}pi\} = 2.40 \textbackslash{}times 10\textasciicircum{}\{-28\} \textbackslash{}text\{ kg\} \textbackslash{}) as shown in figure.

\#\# Question

If all three particles are at rest after the collision, find the initial speed of the protons. Energy is conserved in the collision.

\#\# Answer choices

- A: A: \textbackslash{}( 0.420c \textbackslash{})

- B: B: \textbackslash{}( 0.450c \textbackslash{})

- C: C: \textbackslash{}( 0.270c \textbackslash{})

- D: D: \textbackslash{}( 0.360c \textbackslash{})

\#\# Figure caption (auxiliary; use the image as primary evidence)

The image is a schematic of a particle interaction involving protons and a pion. Here is a description:

1. **Protons**:
   - There are two protons depicted as red circles with a "+" symbol inside each.
   - These protons are labeled "Proton" and each has a mass of \textbackslash{}(1.67 \textbackslash{}times 10\textasciicircum{}\{-27\} \textbackslash{}text\{ kg\}\textbackslash{}).
   - Each proton has an arrow labeled "v" pointing towards the other, indicating movement or velocity towards the center.

2. **Interaction**:
   - The protons approach each other as indicated by the arrows.

3. **Resulting Interaction**:
   - Below the protons, a white arrow points downward to indicate a transformation or interaction result.
   - The resulting image shows the two protons with a green circle labeled "Pion" in between them.
   - The pion has a mass of \textbackslash{}(2.40 \textbackslash{}times 10\textasciicircum{}\{-28\} \textbackslash{}text\{ kg\}\textbackslash{}).

This figure illustrates the concept of protons interacting to produce a pion as part of a subatomic interaction process.

\#\# Dataset metadata

Particle Physics; Multi-Formula Reasoning, Physical Model Grounding Reasoning

\paragraph{Recorded answer/context.}
D: D: \textbackslash{}( 0.360c \textbackslash{})

\paragraph{Figure/image path.}
/root/proposal\_for\_physic/science-mango/phyx\_mini\_run/phyx\_data/test\_image/511.png

\paragraph{Formalization target.}
create a compiling Lean file with sorry bodies at `PhyXMiniProblems/problem\_phyx\_mini\_0511.lean`. The Lean declarations must preserve the physical quantities, dimensions or dimensional roles, figure labels, governing-law hypotheses, and final relation expressed by this problem.
Use Mathlib/Physlib names found through LeanExplore where available. If a physics API is missing, introduce faithful local abstractions rather than scalar placeholder aliases.

\begin{theorem}[Physics formalization target]
\label{thm:physics:phyx_mini_0511:target}
\lean{PhyXMiniProblems.ProblemPhyXMini0511.problem_phyx_mini_0511}
\uses{def:physics:phyx-mini-0511:phyxminiproblems-problemphyxmini0511-speedreadout, def:physics:phyx-mini-0511:phyxminiproblems-problemphyxmini0511-vacuumlightspeedreadout, def:physics:phyx-mini-0511:phyxminiproblems-problemphyxmini0511-protonpioncollisionsetup, def:physics:phyx-mini-0511:phyxminiproblems-problemphyxmini0511-matchesheadonthresholdscenario, def:physics:phyx-mini-0511:phyxminiproblems-problemphyxmini0511-matchessuppliedprotonpionfigure, def:physics:phyx-mini-0511:phyxminiproblems-problemphyxmini0511-matchesproblemmassreadouts, def:physics:phyx-mini-0511:phyxminiproblems-problemphyxmini0511-initialspeedfractionoflight, def:physics:phyx-mini-0511:phyxminiproblems-problemphyxmini0511-hasphysicalcollisionparameters, def:physics:phyx-mini-0511:phyxminiproblems-problemphyxmini0511-satisfiesrelativisticcollisionenergylaws, def:physics:phyx-mini-0511:phyxminiproblems-problemphyxmini0511-answerchoice, def:physics:phyx-mini-0511:phyxminiproblems-problemphyxmini0511-matchesdisplayedspeedchoice}
The assigned autoformalize agent should translate this physics problem into Lean declarations in the covered file, with theorem and lemma proof bodies written as `by sorry`.
\end{theorem}
\begin{proof}
This is an autoformalization task, not a proof task. Produce statements that can later be proved by the physics prover without weakening the physical model.
\end{proof}

% --- Archon declaration topology begin ---
\section{Declaration topology}
\begin{definition}[Mass Quantity]
  \label{def:physics:phyx-mini-0511:phyxminiproblems-problemphyxmini0511-massquantity}
  \lean{PhyXMiniProblems.ProblemPhyXMini0511.MassQuantity}
  A nonnegative physical rest mass, independent of a choice of units
\end{definition}
\begin{proof}
  This is the declared model definition of Mass Quantity.
\end{proof}
\begin{definition}[Speed Quantity]
  \label{def:physics:phyx-mini-0511:phyxminiproblems-problemphyxmini0511-speedquantity}
  \lean{PhyXMiniProblems.ProblemPhyXMini0511.SpeedQuantity}
  A nonnegative physical speed, using Physlib's dimensionful speed type
\end{definition}
\begin{proof}
  This is the declared model definition of Speed Quantity.
\end{proof}
\begin{definition}[Signed Velocity Quantity]
  \label{def:physics:phyx-mini-0511:phyxminiproblems-problemphyxmini0511-signedvelocityquantity}
  \lean{PhyXMiniProblems.ProblemPhyXMini0511.SignedVelocityQuantity}
  A signed velocity component along the horizontal collision axis
\end{definition}
\begin{proof}
  This is the declared model definition of Signed Velocity Quantity.
\end{proof}
\begin{definition}[mass In Kilograms]
  \label{def:physics:phyx-mini-0511:phyxminiproblems-problemphyxmini0511-massinkilograms}
  \lean{PhyXMiniProblems.ProblemPhyXMini0511.massInKilograms}
  \uses{def:physics:phyx-mini-0511:phyxminiproblems-problemphyxmini0511-massquantity}
  Read a physical mass in coherent SI units, hence in kilograms
\end{definition}
\begin{proof}
  This is the declared model definition of mass In Kilograms.
\end{proof}
\begin{definition}[speed Readout]
  \label{def:physics:phyx-mini-0511:phyxminiproblems-problemphyxmini0511-speedreadout}
  \lean{PhyXMiniProblems.ProblemPhyXMini0511.speedReadout}
  \uses{def:physics:phyx-mini-0511:phyxminiproblems-problemphyxmini0511-speedquantity}
  Read a nonnegative speed in selected length and time units
\end{definition}
\begin{proof}
  This is the declared model definition of speed Readout.
\end{proof}
\begin{definition}[signed Velocity Readout]
  \label{def:physics:phyx-mini-0511:phyxminiproblems-problemphyxmini0511-signedvelocityreadout}
  \lean{PhyXMiniProblems.ProblemPhyXMini0511.signedVelocityReadout}
  \uses{def:physics:phyx-mini-0511:phyxminiproblems-problemphyxmini0511-signedvelocityquantity}
  Read a signed velocity component in selected length and time units
\end{definition}
\begin{proof}
  This is the declared model definition of signed Velocity Readout.
\end{proof}
\begin{definition}[speed In Meters Per Second]
  \label{def:physics:phyx-mini-0511:phyxminiproblems-problemphyxmini0511-speedinmeterspersecond}
  \lean{PhyXMiniProblems.ProblemPhyXMini0511.speedInMetersPerSecond}
  \uses{def:physics:phyx-mini-0511:phyxminiproblems-problemphyxmini0511-speedquantity, def:physics:phyx-mini-0511:phyxminiproblems-problemphyxmini0511-speedreadout}
  Meter-per-second readout of a physical speed
\end{definition}
\begin{proof}
  This is the declared model definition of speed In Meters Per Second.
\end{proof}
\begin{definition}[signed Velocity In Meters Per Second]
  \label{def:physics:phyx-mini-0511:phyxminiproblems-problemphyxmini0511-signedvelocityinmeterspersecond}
  \lean{PhyXMiniProblems.ProblemPhyXMini0511.signedVelocityInMetersPerSecond}
  \uses{def:physics:phyx-mini-0511:phyxminiproblems-problemphyxmini0511-signedvelocityquantity, def:physics:phyx-mini-0511:phyxminiproblems-problemphyxmini0511-signedvelocityreadout}
  Meter-per-second readout of a signed velocity component
\end{definition}
\begin{proof}
  This is the declared model definition of signed Velocity In Meters Per Second.
\end{proof}
\begin{definition}[vacuum Light Speed Readout]
  \label{def:physics:phyx-mini-0511:phyxminiproblems-problemphyxmini0511-vacuumlightspeedreadout}
  \lean{PhyXMiniProblems.ProblemPhyXMini0511.vacuumLightSpeedReadout}
  \uses{def:physics:phyx-mini-0511:phyxminiproblems-problemphyxmini0511-signedvelocityreadout}
  Physlib's speed of light has a real-valued dimensionful carrier. Its readout is therefore kept separate from the nonnegative speedReadout above
\end{definition}
\begin{proof}
  This is the declared model definition of vacuum Light Speed Readout.
\end{proof}
\begin{definition}[vacuum Light Speed In Meters Per Second]
  \label{def:physics:phyx-mini-0511:phyxminiproblems-problemphyxmini0511-vacuumlightspeedinmeterspersecond}
  \lean{PhyXMiniProblems.ProblemPhyXMini0511.vacuumLightSpeedInMetersPerSecond}
  \uses{def:physics:phyx-mini-0511:phyxminiproblems-problemphyxmini0511-vacuumlightspeedreadout}
  The coherent SI value of the physical speed of light
\end{definition}
\begin{proof}
  This is the declared model definition of vacuum Light Speed In Meters Per Second.
\end{proof}
\begin{definition}[energy In Joules]
  \label{def:physics:phyx-mini-0511:phyxminiproblems-problemphyxmini0511-energyinjoules}
  \lean{PhyXMiniProblems.ProblemPhyXMini0511.energyInJoules}
  Read a physical energy in coherent SI units, hence in joules
\end{definition}
\begin{proof}
  This is the declared model definition of energy In Joules.
\end{proof}
\begin{definition}[Particle Label]
  \label{def:physics:phyx-mini-0511:phyxminiproblems-problemphyxmini0511-particlelabel}
  \lean{PhyXMiniProblems.ProblemPhyXMini0511.ParticleLabel}
  The three individually tracked particles in the reaction diagram
\end{definition}
\begin{proof}
  This is the declared model definition of Particle Label.
\end{proof}
\begin{definition}[Particle Species]
  \label{def:physics:phyx-mini-0511:phyxminiproblems-problemphyxmini0511-particlespecies}
  \lean{PhyXMiniProblems.ProblemPhyXMini0511.ParticleSpecies}
  The two particle species occurring in the reaction
\end{definition}
\begin{proof}
  This is the declared model definition of Particle Species.
\end{proof}
\begin{definition}[particle Species]
  \label{def:physics:phyx-mini-0511:phyxminiproblems-problemphyxmini0511-particlespecies-value}
  \lean{PhyXMiniProblems.ProblemPhyXMini0511.particleSpecies}
  \uses{def:physics:phyx-mini-0511:phyxminiproblems-problemphyxmini0511-particlelabel, def:physics:phyx-mini-0511:phyxminiproblems-problemphyxmini0511-particlespecies}
  Species attached to each individual particle label
\end{definition}
\begin{proof}
  This is the declared model definition of particle Species.
\end{proof}
\begin{definition}[Electric Charge Class]
  \label{def:physics:phyx-mini-0511:phyxminiproblems-problemphyxmini0511-electricchargeclass}
  \lean{PhyXMiniProblems.ProblemPhyXMini0511.ElectricChargeClass}
  Coarse electric-charge information supplied by the + signs and prose
\end{definition}
\begin{proof}
  This is the declared model definition of Electric Charge Class.
\end{proof}
\begin{definition}[particle Charge Class]
  \label{def:physics:phyx-mini-0511:phyxminiproblems-problemphyxmini0511-particlechargeclass}
  \lean{PhyXMiniProblems.ProblemPhyXMini0511.particleChargeClass}
  \uses{def:physics:phyx-mini-0511:phyxminiproblems-problemphyxmini0511-particlelabel, def:physics:phyx-mini-0511:phyxminiproblems-problemphyxmini0511-electricchargeclass}
  Charge class of the two protons and the neutral pion
\end{definition}
\begin{proof}
  This is the declared model definition of particle Charge Class.
\end{proof}
\begin{definition}[Collision Stage]
  \label{def:physics:phyx-mini-0511:phyxminiproblems-problemphyxmini0511-collisionstage}
  \lean{PhyXMiniProblems.ProblemPhyXMini0511.CollisionStage}
  State immediately before or immediately after the collision
\end{definition}
\begin{proof}
  This is the declared model definition of Collision Stage.
\end{proof}
\begin{definition}[Figure Slot]
  \label{def:physics:phyx-mini-0511:phyxminiproblems-problemphyxmini0511-figureslot}
  \lean{PhyXMiniProblems.ProblemPhyXMini0511.FigureSlot}
  Left, center, and right locations in a row of the supplied diagram
\end{definition}
\begin{proof}
  This is the declared model definition of Figure Slot.
\end{proof}
\begin{definition}[Horizontal Direction]
  \label{def:physics:phyx-mini-0511:phyxminiproblems-problemphyxmini0511-horizontaldirection}
  \lean{PhyXMiniProblems.ProblemPhyXMini0511.HorizontalDirection}
  Horizontal direction of an incoming velocity arrow
\end{definition}
\begin{proof}
  This is the declared model definition of Horizontal Direction.
\end{proof}
\begin{definition}[Figure Color]
  \label{def:physics:phyx-mini-0511:phyxminiproblems-problemphyxmini0511-figurecolor}
  \lean{PhyXMiniProblems.ProblemPhyXMini0511.FigureColor}
  The red and green particle colors used by the supplied figure
\end{definition}
\begin{proof}
  This is the declared model definition of Figure Color.
\end{proof}
\begin{definition}[Proton Pion Collision Figure]
  \label{def:physics:phyx-mini-0511:phyxminiproblems-problemphyxmini0511-protonpioncollisionfigure}
  \lean{PhyXMiniProblems.ProblemPhyXMini0511.ProtonPionCollisionFigure}
  \uses{def:physics:phyx-mini-0511:phyxminiproblems-problemphyxmini0511-particlelabel, def:physics:phyx-mini-0511:phyxminiproblems-problemphyxmini0511-particlespecies, def:physics:phyx-mini-0511:phyxminiproblems-problemphyxmini0511-collisionstage, def:physics:phyx-mini-0511:phyxminiproblems-problemphyxmini0511-figureslot, def:physics:phyx-mini-0511:phyxminiproblems-problemphyxmini0511-horizontaldirection, def:physics:phyx-mini-0511:phyxminiproblems-problemphyxmini0511-figurecolor}
  The label and geometry content of the primary image. It stores particle placement, inward arrows, colors, labels, signs, and printed mass readouts, but deliberately contains no numerical value for the unknown speed v
\end{definition}
\begin{proof}
  This is the declared model definition of Proton Pion Collision Figure.
\end{proof}
\begin{definition}[Proton Pion Collision Setup]
  \label{def:physics:phyx-mini-0511:phyxminiproblems-problemphyxmini0511-protonpioncollisionsetup}
  \lean{PhyXMiniProblems.ProblemPhyXMini0511.ProtonPionCollisionSetup}
  \uses{def:physics:phyx-mini-0511:phyxminiproblems-problemphyxmini0511-massquantity, def:physics:phyx-mini-0511:phyxminiproblems-problemphyxmini0511-speedquantity, def:physics:phyx-mini-0511:phyxminiproblems-problemphyxmini0511-signedvelocityquantity, def:physics:phyx-mini-0511:phyxminiproblems-problemphyxmini0511-particlelabel, def:physics:phyx-mini-0511:phyxminiproblems-problemphyxmini0511-collisionstage, def:physics:phyx-mini-0511:phyxminiproblems-problemphyxmini0511-protonpioncollisionfigure}
  All unknown physical quantities in the collision. Particle energies and total energies are independent dimensionful fields constrained by the laws below. In particular, initialProtonSpeed is not defined from an answer choice or from the desired closed form
\end{definition}
\begin{proof}
  This is the declared model definition of Proton Pion Collision Setup.
\end{proof}
\begin{definition}[particle Rest Mass]
  \label{def:physics:phyx-mini-0511:phyxminiproblems-problemphyxmini0511-particlerestmass}
  \lean{PhyXMiniProblems.ProblemPhyXMini0511.particleRestMass}
  \uses{def:physics:phyx-mini-0511:phyxminiproblems-problemphyxmini0511-massquantity, def:physics:phyx-mini-0511:phyxminiproblems-problemphyxmini0511-particlelabel, def:physics:phyx-mini-0511:phyxminiproblems-problemphyxmini0511-protonpioncollisionsetup}
  Rest mass of a labelled particle in a collision setup
\end{definition}
\begin{proof}
  This is the declared model definition of particle Rest Mass.
\end{proof}
\begin{definition}[Matches Head On Threshold Scenario]
  \label{def:physics:phyx-mini-0511:phyxminiproblems-problemphyxmini0511-matchesheadonthresholdscenario}
  \lean{PhyXMiniProblems.ProblemPhyXMini0511.MatchesHeadOnThresholdScenario}
  \uses{def:physics:phyx-mini-0511:phyxminiproblems-problemphyxmini0511-speedreadout, def:physics:phyx-mini-0511:phyxminiproblems-problemphyxmini0511-signedvelocityreadout, def:physics:phyx-mini-0511:phyxminiproblems-problemphyxmini0511-particlelabel, def:physics:phyx-mini-0511:phyxminiproblems-problemphyxmini0511-protonpioncollisionsetup}
  The particle-content and kinematic facts from the prose. Initially only the two protons are present and their velocity components are equal and opposite. Afterward both protons still exist, the pion has been produced, and all three particles are at rest. None of these facts fixes the magnitude of the initial speed
\end{definition}
\begin{proof}
  This is the declared model definition of Matches Head On Threshold Scenario.
\end{proof}
\begin{definition}[Matches Supplied Proton Pion Figure]
  \label{def:physics:phyx-mini-0511:phyxminiproblems-problemphyxmini0511-matchessuppliedprotonpionfigure}
  \lean{PhyXMiniProblems.ProblemPhyXMini0511.MatchesSuppliedProtonPionFigure}
  \uses{def:physics:phyx-mini-0511:phyxminiproblems-problemphyxmini0511-protonpioncollisionsetup}
  Primary-image evidence: the two inward arrows labelled v, red proton labels and + signs, the downward reaction arrow, and a final green pion between the surviving protons. Printed masses are linked to physical masses separately by MatchesProblemMassReadouts
\end{definition}
\begin{proof}
  This is the declared model definition of Matches Supplied Proton Pion Figure.
\end{proof}
\begin{definition}[Matches Problem Mass Readouts]
  \label{def:physics:phyx-mini-0511:phyxminiproblems-problemphyxmini0511-matchesproblemmassreadouts}
  \lean{PhyXMiniProblems.ProblemPhyXMini0511.MatchesProblemMassReadouts}
  \uses{def:physics:phyx-mini-0511:phyxminiproblems-problemphyxmini0511-massinkilograms, def:physics:phyx-mini-0511:phyxminiproblems-problemphyxmini0511-protonpioncollisionsetup}
  The rest-mass data stated in the problem and printed in the image, expressed as exact rational SI readouts: 1.67 × 10⁻²⁷ kg = 167 / 10²⁹ kg for either proton; 2.40 × 10⁻²⁸ kg = 24 / 10²⁹ kg for the neutral pion. No speed or answer choice occurs in these data
\end{definition}
\begin{proof}
  This is the declared model definition of Matches Problem Mass Readouts.
\end{proof}
\begin{definition}[initial Speed Fraction Of Light]
  \label{def:physics:phyx-mini-0511:phyxminiproblems-problemphyxmini0511-initialspeedfractionoflight}
  \lean{PhyXMiniProblems.ProblemPhyXMini0511.initialSpeedFractionOfLight}
  \uses{def:physics:phyx-mini-0511:phyxminiproblems-problemphyxmini0511-speedinmeterspersecond, def:physics:phyx-mini-0511:phyxminiproblems-problemphyxmini0511-vacuumlightspeedinmeterspersecond, def:physics:phyx-mini-0511:phyxminiproblems-problemphyxmini0511-protonpioncollisionsetup}
  The unknown initial proton speed as a dimensionless fraction of c
\end{definition}
\begin{proof}
  This is the declared model definition of initial Speed Fraction Of Light.
\end{proof}
\begin{definition}[velocity Fraction Of Light]
  \label{def:physics:phyx-mini-0511:phyxminiproblems-problemphyxmini0511-velocityfractionoflight}
  \lean{PhyXMiniProblems.ProblemPhyXMini0511.velocityFractionOfLight}
  \uses{def:physics:phyx-mini-0511:phyxminiproblems-problemphyxmini0511-signedvelocityquantity, def:physics:phyx-mini-0511:phyxminiproblems-problemphyxmini0511-signedvelocityinmeterspersecond, def:physics:phyx-mini-0511:phyxminiproblems-problemphyxmini0511-vacuumlightspeedinmeterspersecond}
  Absolute speed fraction associated with a signed velocity component
\end{definition}
\begin{proof}
  This is the declared model definition of velocity Fraction Of Light.
\end{proof}
\begin{definition}[Has Physical Collision Parameters]
  \label{def:physics:phyx-mini-0511:phyxminiproblems-problemphyxmini0511-hasphysicalcollisionparameters}
  \lean{PhyXMiniProblems.ProblemPhyXMini0511.HasPhysicalCollisionParameters}
  \uses{def:physics:phyx-mini-0511:phyxminiproblems-problemphyxmini0511-massinkilograms, def:physics:phyx-mini-0511:phyxminiproblems-problemphyxmini0511-speedinmeterspersecond, def:physics:phyx-mini-0511:phyxminiproblems-problemphyxmini0511-vacuumlightspeedinmeterspersecond, def:physics:phyx-mini-0511:phyxminiproblems-problemphyxmini0511-energyinjoules, def:physics:phyx-mini-0511:phyxminiproblems-problemphyxmini0511-particlelabel, def:physics:phyx-mini-0511:phyxminiproblems-problemphyxmini0511-collisionstage, def:physics:phyx-mini-0511:phyxminiproblems-problemphyxmini0511-protonpioncollisionsetup, def:physics:phyx-mini-0511:phyxminiproblems-problemphyxmini0511-initialspeedfractionoflight}
  Positivity, subluminality, and nondegeneracy assumptions for the physical collision. They constrain the domain of the relativistic square root but do not specify the requested speed
\end{definition}
\begin{proof}
  This is the declared model definition of Has Physical Collision Parameters.
\end{proof}
\begin{definition}[Satisfies Relativistic Collision Energy Laws]
  \label{def:physics:phyx-mini-0511:phyxminiproblems-problemphyxmini0511-satisfiesrelativisticcollisionenergylaws}
  \lean{PhyXMiniProblems.ProblemPhyXMini0511.SatisfiesRelativisticCollisionEnergyLaws}
  \uses{def:physics:phyx-mini-0511:phyxminiproblems-problemphyxmini0511-massinkilograms, def:physics:phyx-mini-0511:phyxminiproblems-problemphyxmini0511-vacuumlightspeedinmeterspersecond, def:physics:phyx-mini-0511:phyxminiproblems-problemphyxmini0511-energyinjoules, def:physics:phyx-mini-0511:phyxminiproblems-problemphyxmini0511-particlelabel, def:physics:phyx-mini-0511:phyxminiproblems-problemphyxmini0511-collisionstage, def:physics:phyx-mini-0511:phyxminiproblems-problemphyxmini0511-protonpioncollisionsetup, def:physics:phyx-mini-0511:phyxminiproblems-problemphyxmini0511-particlerestmass, def:physics:phyx-mini-0511:phyxminiproblems-problemphyxmini0511-velocityfractionoflight}
  The standard relativistic particle-energy law E = γ(β) m c², additivity of particle energies in each state, and conservation of total energy through the collision. The particle-energy law applies uniformly to every present particle; it neither singles out an answer choice nor contains a solved value of the initial speed
\end{definition}
\begin{proof}
  This is the declared model definition of Satisfies Relativistic Collision Energy Laws.
\end{proof}
\begin{definition}[Answer Choice]
  \label{def:physics:phyx-mini-0511:phyxminiproblems-problemphyxmini0511-answerchoice}
  \lean{PhyXMiniProblems.ProblemPhyXMini0511.AnswerChoice}
  Labels of the four speed fractions printed with the problem
\end{definition}
\begin{proof}
  This is the declared model definition of Answer Choice.
\end{proof}
\begin{definition}[displayed Speed Fraction Of Light]
  \label{def:physics:phyx-mini-0511:phyxminiproblems-problemphyxmini0511-displayedspeedfractionoflight}
  \lean{PhyXMiniProblems.ProblemPhyXMini0511.displayedSpeedFractionOfLight}
  \uses{def:physics:phyx-mini-0511:phyxminiproblems-problemphyxmini0511-answerchoice}
  The coefficient of c printed beside each answer label
\end{definition}
\begin{proof}
  This is the declared model definition of displayed Speed Fraction Of Light.
\end{proof}
\begin{definition}[Matches Displayed Speed Choice]
  \label{def:physics:phyx-mini-0511:phyxminiproblems-problemphyxmini0511-matchesdisplayedspeedchoice}
  \lean{PhyXMiniProblems.ProblemPhyXMini0511.MatchesDisplayedSpeedChoice}
  \uses{def:physics:phyx-mini-0511:phyxminiproblems-problemphyxmini0511-protonpioncollisionsetup, def:physics:phyx-mini-0511:phyxminiproblems-problemphyxmini0511-initialspeedfractionoflight, def:physics:phyx-mini-0511:phyxminiproblems-problemphyxmini0511-answerchoice, def:physics:phyx-mini-0511:phyxminiproblems-problemphyxmini0511-displayedspeedfractionoflight}
  A speed ratio agrees with a value displayed to three decimal places when its absolute error is less than half of 0.001
\end{definition}
\begin{proof}
  This is the declared model definition of Matches Displayed Speed Choice.
\end{proof}
\begin{lemma}[initial Proton Lorentz Factor eq 179 div 167]
  \label{lem:physics:phyx-mini-0511:phyxminiproblems-problemphyxmini0511-initialprotonlorentzfactor-eq-179-div-167}
  \lean{PhyXMiniProblems.ProblemPhyXMini0511.initialProtonLorentzFactor_eq_179_div_167}
  \uses{def:physics:phyx-mini-0511:phyxminiproblems-problemphyxmini0511-protonpioncollisionsetup, def:physics:phyx-mini-0511:phyxminiproblems-problemphyxmini0511-matchesheadonthresholdscenario, def:physics:phyx-mini-0511:phyxminiproblems-problemphyxmini0511-matchesproblemmassreadouts, def:physics:phyx-mini-0511:phyxminiproblems-problemphyxmini0511-initialspeedfractionoflight, def:physics:phyx-mini-0511:phyxminiproblems-problemphyxmini0511-hasphysicalcollisionparameters, def:physics:phyx-mini-0511:phyxminiproblems-problemphyxmini0511-satisfiesrelativisticcollisionenergylaws}
  At threshold, energy conservation first yields γ = (2 M\_p + M\_π) / (2 M\_p) = 179 / 167. This is an intermediate consequence of the governing laws and mass data, not an assumed field of the collision setup
\end{lemma}
\begin{proof}
  The conclusion follows from the governing relations and figure readouts named in the statement of initial Proton Lorentz Factor eq 179 div 167.
\end{proof}
% --- Archon declaration topology end ---
% --- Archon physics formalization source end ---
